\section{Reproducibility and Artifact Identity}

The release record binds source, containers, model material, workload, raw
measurements, telemetry, and rendered reports. Local Docker image IDs and
unpacked sizes are distinct from registry manifest digests and compressed
sizes; neither is inferred from the other. Native and stock model-artifact
inventories are also recorded independently even though both share the same
repository revision.

\begin{table}[t]
  \centering
  \caption{Digest-bound artifact identity. Every pending value must be replaced
  from the accepted production manifest before publication.}
  \label{tab:artifact-identity}
  \small
  \begin{tabularx}{\textwidth}{>{\raggedright\arraybackslash}p{0.31\textwidth}X}
    \toprule
    Record & Value \\
    \midrule
    Evidence state & \PaperEvidenceState \\
    Native Git commit & \NativeCommit \\
    Native local image ID & \NativeLocalImageId \\
    Native OCI digest & \NativeImageDigest \\
    Stock local image ID & \StockLocalImageId \\
    Stock OCI digest & \StockImageDigest \\
    Shared model revision & \code{5ecdb67327fd37bb2e042aab12ff7391903235d3} \\
    Native model-artifact SHA-256 & \NativeModelArtifactSha \\
    Stock model-artifact SHA-256 & \StockModelArtifactSha \\
    Evidence-manifest SHA-256 & \EvidenceManifestSha \\
    Evidence timestamp & \EvidenceTimestamp \\
    Evidence bundle & \EvidenceBundleUri \\
    \bottomrule
  \end{tabularx}
\end{table}

The intended public code location is
\url{https://github.com/luka-loehr/qwen3-tts-native}. At publication time, the
paper must point to an immutable release tag and OCI digest, not only to a
mutable branch. The release archive must contain the English LaTeX source,
bibliography, vector figures, generated plot data, final paper PDF, benchmark
report PDF, evidence manifest, and checksums. Large raw evidence may be hosted
separately when the repository cannot carry it, but the public URI and manifest
digest must remain stable.

Reproduction requires a clean pull by immutable image digest, an idle supported
GPU, the documented hardened run command, a successful full-pipeline readiness
check, and the external workload client. Reproducers should validate the
manifest before generating tables or plots and should report any hardware,
driver, power-mode, workload, or comparator difference rather than presenting
it as the original experiment.
